\documentclass[11pt]{article}
\usepackage{amsmath,amssymb,amsthm,booktabs,geometry,graphicx}
\usepackage[utf8]{inputenc}
\usepackage[T1]{fontenc}
\usepackage[round,authoryear]{natbib}
\usepackage[colorlinks=true,linkcolor=black,citecolor=black,urlcolor=blue!50!black]{hyperref}
\geometry{margin=1in}
\setlength{\emergencystretch}{2em}

\newtheorem{proposition}{Proposition}
\newtheorem{corollary}{Corollary}
\newtheorem{assumption}{Assumption}
\theoremstyle{definition}
\newtheorem{postulate}{Postulate}
\newtheorem{definition}{Definition}
\theoremstyle{remark}
\newtheorem*{remark}{Remark}

\newcommand{\JJ}{\mathcal{J}}
\newcommand{\DD}{\mathcal{D}}
\newcommand{\Dun}{\mathcal{D}_1}
\newcommand{\KL}[2]{D_{KL}\!\left(#1\,\Vert\,#2\right)}

\title{Theoretical Supplement\\[4pt]
\large Information-theoretic foundations of an entropic systemic-stress\\indicator: what the main paper leaves out}
\author{Supplement to ``Entropy Jumps and Information Relaxation\\in Financial Markets''}
\date{October 2026}

\begin{document}
\maketitle

\begin{abstract}
\noindent The main paper is refocused on a single object aimed at central banks:
the dependence channel $-\tfrac12\log\det R$ of a multi-market panel, its exact
decomposition into within-block and between-block terms, and the pricing in
nats of published stress scenarios. This document collects everything that
refocusing removes from the body of the paper, to serve as an appendix or to be
supplied on request. It contains the bridge from maximum entropy in volume time
to geometric Brownian motion, the decomposition of entropy into three channels
and the structural index $\JJ=\KL{p}{\prod_i q_i}$, the jump-diffusion dynamics
(one theorem, two postulates, and the Ornstein--Uhlenbeck structure forced by
Doob's theorem), the empirical programme H1--H8 with its verdict against the
calibrated nulls, and the full geometry of Kullback--Leibler-ball scenarios. Two
results are new relative to arXiv~v1. The first is a block decomposition
identity for the dependence channel, which links this document to the main
paper. The second shows that the full-dimensional lift of a portfolio scenario
that minimises the ambient divergence costs exactly the radius $\eta$ and has an
asset-by-asset reading (the shock is carried by each asset's covariance with
the portfolio), while the default uniform dilation costs $N$~times as much in
variance. The second point corrects a remark in v1.
\end{abstract}

\tableofcontents

\section*{Foreword: what this document contains, and why}
\addcontentsline{toc}{section}{Foreword}

\paragraph{How it relates to the main paper.} The October 2026 assessment found
that a single result of arXiv~v1 clears both calibrated nulls on both panels:
the regime gap of the dependence channel $H^{dep}=\tfrac12\log\det R$. Phase~A
confirmed it in three forms (raw, Forbes--Rigobon devolatilised, rank-based) and
showed that the absorption ratio of \citet{kritzman2011}, the benchmark measure
of the same family, fails the block null on both panels. The main paper is
therefore rebuilt around an ``entropic CISS'' in the sense of
\citet{hollo2012}: the dependence channel of a multi-market panel, its
within/between-block decomposition, central-bank-style validation, and the
pricing of published scenarios in nats.

Everything else leaves the body of the paper. Not because it is wrong: the
mathematical results remain exact. It is because these elements are not needed
for the indicator, or because they do not survive the nulls and therefore
cannot ground a surveillance tool. Collecting them here serves two purposes. A
referee who asks where the dependence channel comes from, and why it is this
channel rather than the full structural index, finds the answer. A reader who
wants to take up the theory of information dynamics finds the whole chain, with
the exact status of each link.

\paragraph{Status of each statement.} The document separates three categories
and flags them throughout: \emph{theorems} (exact under their assumptions),
\emph{postulates} (substantive, testable economic assumptions), and
\emph{empirical findings} (with the null they clear or fail).
Table~\ref{tab:statut} gives the map.

\begin{table}[htbp]
\centering
\small
\begin{tabular}{@{}p{5.0cm}p{3.0cm}p{4.6cm}l@{}}
\toprule
Item & Status & Empirical verdict & Section \\
\midrule
Gaussian maximum entropy & theorem & --- & \ref{sec:bridge} \\
Volume bridge (MDH, subordinated Brownian motion) & argument under assumptions & not tested here & \ref{sec:bridge} \\
Three-channel decomposition & exact identity & --- & \ref{sec:channels} \\
Structural index $\JJ$ as a divergence & exact identity & range clears both nulls & \ref{sec:channels} \\
Downward entropy jumps & theorem (on average) & not directly estimable & \ref{sec:dynamics} \\
Relaxation (P1) & postulate & half-life not identified & \ref{sec:dynamics} \\
One-sided constraint (P2) & postulate & not established (H5) & \ref{sec:dynamics} \\
Tail--correlation coupling & consequence & not identified (H6) & \ref{sec:dynamics} \\
Dependence channel under stress & finding & \textbf{clears both nulls} & \ref{sec:program}, \ref{sec:blocks} \\
KL ball, $1$~nat $=\sqrt2\,\sigma$ & theorem & --- & \ref{sec:stress} \\
Minimal lift (rank one) & theorem (new) & --- & \ref{sec:stress} \\
Radius calibrated on the flow & rule & too mild against history & \ref{sec:stress} \\
Block decomposition & identity (new here) & subject of the main paper & \ref{sec:blocks} \\
\bottomrule
\end{tabular}
\caption{Map of the statements in this document and their status.}
\label{tab:statut}
\end{table}

\paragraph{Provenance.} Sections~\ref{sec:bridge} to \ref{sec:dynamics} and
Appendices~\ref{app:units} to \ref{app:mi} take up, with updated
cross-references, the text of the reference French version
(\texttt{paper/main\_fr.tex}), which supersedes \texttt{paper/main.tex} for
the dynamics section and the appendices. Sections~\ref{sec:program} to
\ref{sec:open} are written for this document. Every number quoted comes from the
same code (\texttt{python -m src.run\_empirical}, two-hundred-replication nulls,
Phase~A); none was recomputed for this supplement.

\section*{Notation}
\addcontentsline{toc}{section}{Notation}

\noindent Two warnings before the table. First, differential entropy $h$
\textbf{depends on units}: $h(aX)=h(X)+\log|a|$, so expressing returns in
percent rather than decimals shifts $h$ by a constant. Divergences
($\KL{\cdot}{\cdot}$, $I$, $\DD$, $\JJ$), by contrast, are invariant under any
invertible change of variable: that is the reason for the structural index.
Second, $K$ denotes kurtosis and ``KL'' stands for Kullback and Leibler; the two
are unrelated.

\begin{center}
\small
\begin{tabular}{@{}lll@{}}
\toprule
Symbol & Meaning & Nature \\
\midrule
$p_t$ & true conditional return distribution at $t$ & --- \\
$q_t$ & Gaussian with the same mean and covariance as $p_t$ & --- \\
$h(\cdot)$ & differential entropy & nats, \emph{unit-dependent} \\
$\KL{p}{q}$ & Kullback--Leibler divergence & nats, invariant \\
$\DD_t$ & entropy deficit $=\KL{p_t}{q_t}\ge0$ & nats, invariant \\
$\Dun(z)$ & deficit of a standardised scalar (negentropy) & nats, invariant \\
$H^{vol}_t$ & scale channel $=\sum_i\log\sigma_{i,t}$ & nats, \emph{unit-dependent} \\
$H^{dep}_t$ & dependence channel $=\tfrac12\log\det R_t\le0$ & nats, invariant \\
$H^{cov}_t$ & $H^{vol}_t+H^{dep}_t$; the constant $\tfrac N2\log2\pi e$ is omitted & nats \\
$\JJ_t$ & structural index $=\DD_t-H^{dep}_t\ge0$ & nats, invariant \\
$I(X;Y)$ & mutual information & nats, invariant \\
$I^{(h)}_t$ & information flow $=\KL{p_t}{p_{t-h}}$ & nats \\
$S$, $K$ & skewness, excess kurtosis & dimensionless \\
$\Sigma_t=D_tR_tD_t$ & covariance, standard deviations, correlation & --- \\
$\lambda_k$, $N^{\text{eff}}_t$ & eigenvalues of $R_t$, effective number of modes & $1\le N^{\text{eff}}\le N$ \\
$\kappa$, $\sigma_\JJ$ & relaxation rate and diffusion coefficient of $\JJ$ & --- \\
$\eta$ & radius of a stress scenario & nats \\
$N$ & number of assets & --- \\
\bottomrule
\end{tabular}
\end{center}

\section{From order flow to a Gaussian price: the volume bridge}
\label{sec:bridge}

\subsection{The price as accumulated order flow}

Over $[t,t+\Delta]$ let the market execute $M$ trades, trade $k$ carrying signed
volume $\epsilon_k v_k$ with $\epsilon_k\in\{-1,+1\}$ the sign (buyer- or
seller-initiated) and $v_k>0$ the size. Write $V=\sum_k v_k$ for total volume
and $Q=\sum_k \epsilon_k v_k$ for net signed order flow.

\begin{assumption}[Linear impact]\label{as:impact}
The log-price responds linearly to net signed order flow,
$\Delta X_{t,\Delta} = X_{t+\Delta}-X_t = \lambda\,Q$, with $\lambda>0$ constant
over the interval.
\end{assumption}

This is Kyle's \citeyearpar{kyle1985} equilibrium: the market maker, unable to
distinguish informed from uninformed flow, sets a linear pricing rule and
$\lambda$ is the inverse depth. It is the definition of ``price as a function of
traded volume'' that the rest of the section needs, and it is the weakest link
in the chain; Section~\ref{sec:bridge-caveats} says what happens when it fails.

\begin{assumption}[Exchangeable signs and a volume-proportional variance]
\label{as:signs}
Conditional on $V$ and on the information available at $t$, the market's belief
about $Q$ is constrained only by its first two moments, with
$\mathbb{E}[Q\mid V] = m$ and $\mathrm{Var}(Q\mid V) = \varsigma^2 V$.
\end{assumption}

The variance form is what one gets when trade signs are exchangeable with
bounded correlation and sizes have a stable mean: for independent signs,
$\mathrm{Var}(Q)=\sum_k v_k^2 \approx \bar v V$. The economic content is that
\emph{volume, not time, is the natural measure of how much can happen}.

\subsection{Maximum entropy in volume time}

\begin{proposition}[Maximum entropy under first and second moments]
\label{prop:maxent}
Among absolutely continuous distributions on $\mathbb{R}^N$ with mean $\mu$ and
covariance $\Sigma\succ0$, the Gaussian $\mathcal N(\mu,\Sigma)$ uniquely
maximises differential entropy, and
$h^{\max}(\Sigma)=\tfrac12\log\left((2\pi e)^N\det\Sigma\right)$.
\end{proposition}

\begin{proof}
Let $q=\mathcal N(\mu,\Sigma)$ and let $p$ have the same first two moments and
finite differential entropy. Since $\log q$ is an exact quadratic, its
expectation depends only on those moments, so
$\mathbb{E}_p[\log q]=\mathbb{E}_q[\log q]=-h(q)$. Non-negativity of relative
entropy gives
$0\le \KL{p}{q} = -h(p)+h(q)$, hence $h(p)\le h(q)$, with equality iff $p=q$
almost everywhere. The argument requires no differentiability of $p$ and no
prior existence of a maximiser.
\end{proof}

The variational route -- a Lagrangian in the normalisation, mean and covariance
constraints, whose first-order condition makes $\log p$ quadratic -- explains
\emph{why} the maximiser is Gaussian, but the divergence argument is what
establishes optimality and uniqueness.

Applying Proposition~\ref{prop:maxent} to Assumption~\ref{as:signs} in one
dimension and mapping through Assumption~\ref{as:impact}:

\begin{equation}
\label{eq:mdh}
\Delta X_{t,\Delta}\ \big|\ V \ \sim\ \mathcal{N}\!\left(\lambda m,\ \lambda^2\varsigma^2 V\right).
\end{equation}

Equation~\eqref{eq:mdh} is the bridge. Three things about it deserve emphasis.

First, \emph{the conditioning variable is volume, not time}. Maximum entropy
does not deliver a variance; it delivers the least-committal distribution given
one. Assumption~\ref{as:signs} supplies the variance from an observable, and
that observable is traded volume. Equation~\eqref{eq:mdh} is the
mixture-of-distributions hypothesis of \citet{clark1973} and
\citet{tauchen1983}, obtained here as a maximum-entropy statement rather than
posited; \citet{ane2000} show empirically that returns measured in trade time
are close to Gaussian even though calendar-time returns are not, which is the
prediction \eqref{eq:mdh} makes.

Second, \emph{Gaussianity is conditional}. Unconditionally,
$\Delta X$ is a scale mixture of normals over the random volume $V$, hence
leptokurtic even when \eqref{eq:mdh} holds exactly. This is the
\citet{mandelbrot1967} subordination mechanism, and it means fat tails are not
by themselves evidence against maximum entropy.

Third, the multivariate extension is immediate: with $N$ assets,
$\Delta \mathbf{X}\mid V \sim \mathcal N(\boldsymbol\mu V, \Sigma V)$, and
Proposition~\ref{prop:maxent} applies as stated.

\subsection{From Gaussian increments to Brownian motion and to GBM}

Let $\theta_t=\int_0^t v_s\,ds$ be the cumulative volume clock. Equation
\eqref{eq:mdh} gives the one-increment distribution; a process needs more.

\begin{assumption}[Temporal consistency]\label{as:time}
In the volume clock, increments over disjoint intervals are independent, sample
paths are continuous, and the covariance of an increment is proportional to
elapsed volume.
\end{assumption}

Under Assumption~\ref{as:time} with constant coefficients there is a matrix $B$
with $BB^{\!\top}=\Sigma$ such that $X_\theta = X_0 + \mu\theta + BW_\theta$ in
volume time, so in calendar time
\begin{equation}
\label{eq:subordination}
X_t = X_0 + \mu\,\theta_t + B\,W_{\theta_t},
\end{equation}
a Brownian motion \emph{subordinated} to the volume clock. If trading intensity
is constant, $\theta_t=\bar v\,t$ and \eqref{eq:subordination} is ordinary
Brownian motion in calendar time; exponentiating and applying Itô,
\begin{equation}
\frac{dS_{i,t}}{S_{i,t}} = \left(a_{i,t}+\tfrac12\sigma_{i,t}^2\right)dt + \sigma_{i,t}\,dW_{i,t},
\end{equation}
so the drift of the price is not the drift of the log-price -- the same
second-order correction that will reappear in Section~\ref{sec:channels} when we
differentiate $\log\det R_t$.

It is worth being precise about what Assumption~\ref{as:time} does. By Monroe's
\citeyearpar{monroe1978} theorem, \emph{every} continuous local martingale is a
time-changed Brownian motion, so subordination alone has no empirical content.
The content is in the clock: \eqref{eq:mdh} says the right clock is volume, and
Assumption~\ref{as:time} says the market has no memory in that clock. Maximum
entropy is what selects Brownian motion from the class Monroe's theorem allows.

\subsection{The contrapositive}
\label{sec:contrapositive}

Proposition~\ref{prop:maxent} runs both ways. Write $q_t$ for the Gaussian with
the conditional mean and covariance of $p_t$. Then

\begin{equation}
\label{eq:deficit}
\DD_t \;\equiv\; h(q_t) - h(p_t) \;=\; \KL{p_t}{q_t} \;\ge\; 0,
\end{equation}
with equality if and only if $p_t$ is Gaussian. So: \emph{prices are non-Gaussian
exactly to the extent that their conditional distribution is not
entropy-maximal}, and $\DD_t$ measures the extent. In the language of
Section~\ref{sec:bridge}, a positive $\DD_t$ means constraints beyond the first
two moments are binding -- that somebody knows something the second moment does
not encode.

Two sources produce a positive deficit, and separating them matters. A random
volume clock produces a scale mixture, which is symmetric and leptokurtic.
Binding directional information produces asymmetry. Section~\ref{sec:channels}
splits $\DD$ along exactly this line.

\subsection{Where the bridge is weak}
\label{sec:bridge-caveats}

Assumption~\ref{as:impact} is linear; the empirical impact of a metaorder is
closer to square-root in size, which makes the mapping from flow to price
concave and the induced return distribution thinner-tailed than the flow.
Assumption~\ref{as:signs} treats signs as exchangeable; order flow is in fact
strongly autocorrelated, which inflates $\mathrm{Var}(Q)$ above $\varsigma^2 V$
and makes the volume clock an imperfect variance proxy. Neither failure changes
the \emph{form} of the argument -- maximum entropy given a variance budget is
still Gaussian -- but both mean the volume clock is an approximation to the
information clock rather than the thing itself. We use volume-time reasoning to
motivate the framework and calendar-time estimation throughout the empirical
work, so no result below depends on \eqref{eq:mdh} holding exactly.

\section{Three channels of entropy}
\label{sec:channels}

\subsection{Scale and dependence}

Write $\Sigma_t = D_tR_tD_t$ with $D_t=\mathrm{diag}(\sigma_{1,t},\dots,\sigma_{N,t})$.
Then $\det\Sigma_t = \left(\prod_i\sigma_{i,t}^2\right)\det R_t$ and
\begin{equation}
\label{eq:ceiling}
h(q_t) \;=\; \underbrace{\tfrac{N}{2}\log(2\pi e)}_{\text{constant}}
\;+\; \underbrace{\sum_i \log\sigma_{i,t}}_{H^{vol}_t}
\;+\; \underbrace{\tfrac12\log\det R_t}_{H^{dep}_t}.
\end{equation}
By Hadamard's inequality $H^{dep}_t\le 0$, with equality iff $R_t=I$:
\emph{correlation always destroys entropy}, and $-H^{dep}_t$ is precisely the
Gaussian multi-information (total correlation) of the return vector. A useful
reading of the same quantity is spectral: with $\lambda_k$ the eigenvalues of
$R_t$ and $\tilde\lambda_k=\lambda_k/N$, the participation number
$N^{\text{eff}}_t=\exp\!\left(-\sum_k\tilde\lambda_k\log\tilde\lambda_k\right)$
is the effective number of independent risk modes; it equals $N$ for $R=I$ and
falls toward $1$ as one collective mode absorbs the variance.

Equation~\eqref{eq:ceiling} is the ceiling, not the entropy. It is what the
covariance matrix alone implies, and Proposition~\ref{prop:maxent} says the
truth sits below it.

\subsection{Shape: the entropy deficit and what it is made of}

Subtracting \eqref{eq:deficit} from \eqref{eq:ceiling} gives the paper's
decomposition, three channels and a constant:
\begin{equation}
\label{eq:three}
h(p_t) \;=\; \tfrac{N}{2}\log(2\pi e) \;+\; \underbrace{H^{vol}_t}_{\text{scale}} \;+\; \underbrace{H^{dep}_t}_{\text{dependence}} \;-\; \underbrace{\DD_t}_{\text{shape}}.
\end{equation}

The shape channel is skewness and fat tails, and the Edgeworth expansion says so
explicitly: for a standardised variable with skewness $S$ and excess kurtosis
$K$ near Gaussianity,
\begin{equation}
\label{eq:edgeworth}
J \;\approx\; \frac{S^2}{12} + \frac{K^2}{48}.
\end{equation}
Equation~\eqref{eq:edgeworth} is the answer to the question ``where would an
entropy jump be visible?'': in the third and fourth moments, additively, with an
asymmetry channel and a tail channel. It is also the source of the sign problem
that Section~\ref{sec:dynamics} has to solve, because $S$ enters squared --
\emph{negentropy cannot tell a crash from a melt-up}. As an estimator
\eqref{eq:edgeworth} is unusable on daily returns; Section~\ref{sec:program}
documents how badly and what we use instead.

\subsection{The structural index}

The scale channel is a nuisance for our purposes: it dominates the level of
\eqref{eq:three} numerically, it moves with volatility, and a change in units
changes it. The other two channels are scale-free. Combining them, and using the
exact chain rule
$\KL{p}{\prod_i q_i} = \sum_i \KL{p_i}{q_i} + I(x)$ with $I$ the
multi-information:

\begin{definition}[Structural information index]
\label{def:J}
\begin{equation}
\JJ_t \;\equiv\; \DD_t \;-\; H^{dep}_t \;=\; \KL{p_t}{\textstyle\prod_i q_{i,t}} \;\ge\; 0 ,
\end{equation}
the divergence of the conditional joint return distribution from a product of
Gaussian marginals with the same scales.
\end{definition}

$\JJ$ is zero if and only if returns are independent and Gaussian; it rises when
assets couple, when marginals fatten, or when they skew. It is measured in nats,
it is invariant to the level of volatility, and it is additive across its two
channels, so one can always ask how much of a move came from dependence and how
much from shape. It is the state variable of Section~\ref{sec:dynamics}.


\section{Information dynamics: one theorem, two postulates}
\label{sec:dynamics}

\subsection{Downward jumps are a theorem}

Two distinct statements underpin the downward half of the postulated dynamics.
Conflating them is an easy mistake, and earlier versions of this work made it;
we therefore separate them explicitly.

Recall first the definition of the \emph{mutual information} between the return
vector $X$ and a signal $Y$:
\begin{equation}
\label{eq:mi}
I(X;Y) \;=\; \KL{p_{XY}}{p_X\otimes p_Y}
\;=\; \iint p(x,y)\,\log\frac{p(x,y)}{p(x)p(y)}\,dx\,dy ,
\end{equation}
that is, the divergence of the joint law from the product of the marginals: how
far the joint law is from independence. It is non-negative, symmetric, zero if
and only if $X\perp Y$, and, unlike $h$, invariant under invertible
reparametrisation. For a scalar Gaussian pair with correlation $\rho$,
$I=-\tfrac12\log(1-\rho^2)$.

\begin{proposition}[Fall in average conditional entropy]
\label{prop:conditioning}
Let $X$ be the return vector and $Y$ a signal. Then
\begin{equation*}
\mathbb{E}_Y\!\left[h(X\mid Y)\right] \;=\; h(X) - I(X;Y) \;\le\; h(X),
\end{equation*}
with equality if and only if $X\perp Y$.
\end{proposition}

\begin{proposition}[Monotonicity of the ceiling under constraint]
\label{prop:constraint}
Let $\mathcal{C}\subseteq\mathcal{C}'$ be two constraint sets and
$\mathcal{Q}(\cdot)$ the set of laws that satisfy them. Then
$\mathcal{Q}(\mathcal{C}')\subseteq\mathcal{Q}(\mathcal{C})$ and hence
\begin{equation*}
\sup_{p\,\in\,\mathcal{Q}(\mathcal{C}')} h(p) \;\le\; \sup_{p\,\in\,\mathcal{Q}(\mathcal{C})} h(p).
\end{equation*}
\end{proposition}

Both results are classical \citep{cover2006} and they are not equivalent; the
detailed proof and a worked Gaussian example are in Appendix~\ref{app:mi}. Their
scopes differ along two axes.

Proposition~\ref{prop:conditioning} \emph{quantifies} the fall (exactly
$I(X;Y)$ nats) but only \emph{on average} over realisations of the signal. For a
particular $y$, $h(X\mid Y=y)$ can exceed $h(X)$: a sufficiently confusing signal
leaves the market more uncertain than before. One should therefore expect a jump
component that is mainly, not exclusively, one-sided, which is what
Section~\ref{sec:program} reports.

Proposition~\ref{prop:constraint}, by contrast, holds \emph{for every
realisation}, but it concerns the maximum-entropy \emph{ceiling}, not the actual
entropy, and it quantifies nothing: it says that the ceiling cannot rise, not by
how much it falls.

What the two say together, and what motivates everything that follows: the fall
is \emph{discontinuous}, because the signal arrives at a point in time, and its
expected size in nats is the information content of the news, an amplitude
measured on a scale that depends on neither the asset, the currency nor the
units.

\paragraph{An honesty caveat.} Neither proposition is used to compute anything in
the empirical part. We measure $\JJ$, not $h(X)-\mathbb{E}_Y[h(X\mid Y)]$:
estimating a differential entropy in dimension~20 on a 252-day window is out of
reach (Appendix~\ref{app:units}). The link between the theorem and the
statistics reported later is a line of reasoning, not an estimated identity, and
the reader should treat it as such.

\subsection{Two postulates}

\begin{postulate}[Relaxation]
\label{p:relax}
In the absence of new information, the structural index decays towards a
long-run level at a constant rate:
$d\JJ_t = -\kappa\left(\JJ_t-\bar\JJ\right)dt + \sigma_\JJ\,dW_t$ with $\kappa>0$.
\end{postulate}

The economics behind P\ref{p:relax} is that a constraint goes stale. The news is
priced in, disagreement re-forms, positions disperse, and the cross-section
recovers the freedom the shock took away. Nothing in information theory delivers
this; it is the substantive assumption of the paper, and $\kappa$ is the object
to be estimated.

\paragraph{Why an Ornstein--Uhlenbeck process and not something else?} The choice
is not one of convenience: it is \emph{forced} once three minimal requirements
are imposed, and it is the only process that meets all three.

\begin{enumerate}
\item \emph{Markov.} The future of $\JJ$ depends on the past only through its
current level. This is the parsimony requirement: postulate no memory beyond
what the current state already holds.
\item \emph{Stationary.} $\JJ$ has a finite long-run level and variance. Without
this the index drifts, it is not a state variable, and the regime comparisons of
Section~\ref{sec:program} are meaningless.
\item \emph{Gaussian noise.} Between two information arrivals we claim to know
nothing about the noise beyond its variance: this is the maximum-entropy
principle applied to the noise itself, exactly as it was applied to returns in
Section~\ref{sec:bridge}.
\end{enumerate}

\citet{doob1942}'s theorem states that \textbf{the Ornstein--Uhlenbeck process is
the only continuous-time process that is at once Markov, Gaussian and
stationary}. The three requirements therefore leave no latitude on the diffusive
part: it is determined. The jump term is added precisely to \emph{break}
Gaussianity where the framework says it fails, namely at information arrivals.
The structure of \eqref{eq:ou} is thus an imposed skeleton plus a deliberately
non-Gaussian component, which is exactly the class introduced by
\citet{bns2001}.

\paragraph{What remains arbitrary.} Three things, which must be stated rather
than presented as if they followed from the theorem. The rate $\kappa$ is
assumed \emph{constant}: nothing rules out a crisis being absorbed faster, or
slower, than ordinary news, and the data used here cannot tell. The pull-back is
\emph{linear} in the deviation from the long-run level, whereas a pull-back that
accelerates far from the mean would be just as plausible. Finally, $\JJ$ is
non-negative by construction, whereas an Ornstein--Uhlenbeck process with a
Brownian term can turn negative: the \citet{bns2001} specification avoids this
by dropping the Brownian term and driving the process with a subordinator with
positive jumps, which guarantees a stationary law supported on $\mathbb{R}^+$.
Empirically the issue does not arise ($\JJ$ stays between $2.17$ and
$21.13$~nats on the main panel and between $0.42$ and $1.69$ on the replication
panel, far from the boundary), but formally one should either drop
$\sigma_\JJ$ or write the process on $\log\JJ$.

Nor is the postulate innocent as a matter of measurement: an estimator with
memory will manufacture decay even with no economic relaxation at all, which is
why Section~\ref{sec:program} measures the decay that a memoryless world already
produces.

\begin{postulate}[Asymmetry of the constraint]
\label{p:sign}
The constraint imposed by information is one-sided in the direction of losses,
so that an entropy jump comes with negative skewness of the market factor.
\end{postulate}

\paragraph{Why this postulate is logically necessary.} The index $\JJ$ is an
\emph{even} function of skewness: flipping the sign of skewness leaves its value
unchanged. The reason is immediate. The shape channel is
$\Dun(z)=\KL{p_z}{\varphi}$, and the normal law $\varphi$ is symmetric; so
replacing $z$ by $-z$ leaves the reference law unchanged, and the divergence too:
\begin{equation}
\label{eq:even}
\Dun(-z)=\Dun(z),\qquad\text{and therefore}\qquad \JJ(-x)=\JJ(x).
\end{equation}
The Edgeworth expansion \eqref{eq:edgeworth} makes it visible at a glance:
skewness $S$ enters squared.

The consequence is structural, not anecdotal. \textbf{A distribution with a long
left tail (a crash) and its mirror image with a long right tail (a melt-up) have
exactly the same $\JJ$.} The index measures the \emph{size} of the information,
never its direction. No refinement of the estimator will change this: it is a
property of the object measured. If the paper is to say anything about direction,
it must therefore add an external ingredient. P\ref{p:sign} is that ingredient,
and it is the smallest possible: a single statement, about the sign, added to a
framework that by construction produces none.

\paragraph{What it rests on.} Not on a mechanism we would assert, but on an
established empirical regularity: the \emph{leverage effect} in the sense of
\citet{black1976}, namely that negative returns are followed by larger rises in
volatility than positive returns of the same size. It is one of the most robust
stylised facts in empirical finance \citep{christie1982,engle1993,nelson1991}.
Translated into the language of this paper: for an equal quantity of
information, the market's response is larger on the way down than on the way up,
so a jump in $\JJ$ comes with negative market skewness.

We deliberately abstain from choosing among the competing causal explanations of
this effect (financial leverage proper, volatility feedback, or asymmetric
withdrawal of liquidity on the sell side). The postulate depends on none of them:
it asserts only the asymmetry, not its cause. It is separately testable (H5) and
separately refutable, and Section~\ref{sec:program} concludes that it is not
established.

\subsection{The resulting dynamics}

Proposition~\ref{prop:conditioning} together with P\ref{p:relax} gives a
non-Gaussian Ornstein--Uhlenbeck process with one-sided jumps, of the class
introduced by \citet{bns2001} for stochastic volatility:
\begin{equation}
\label{eq:ou}
d\JJ_t \;=\; -\kappa\left(\JJ_t-\bar\JJ\right)dt \;+\; \sigma_\JJ\,dW_t \;+\; \int_{z>0} z\,\mathcal{P}(dt,dz).
\end{equation}
The integral term needs careful reading, because its notation is compact.

\paragraph{What $\mathcal{P}$ is not.} It is neither a Gaussian law nor a
function. The letter $\mathcal N$ is reserved in this document for the normal
law, so the jump measure carries a distinct symbol.

\paragraph{What $\mathcal{P}$ is.} A \emph{random counting measure}, known as a
Poisson random measure. One can picture it as a random scatter of points in the
(date, size) plane: each point is an information event, its abscissa is the date
it occurs and its ordinate the size of the jump it causes. Formally,
$\mathcal{P}(dt,dz)$ equals $1$ if there is a jump in the time interval
$[t,t+dt)$ whose size falls in $[z,z+dz)$, and $0$ otherwise. The integral
$\int_{z>0} z\,\mathcal{P}(dt,dz)$ then reads simply: ``if there is a jump at
date $t$, add its size $z$ to $\JJ$; otherwise add nothing''. It is a compact
way of writing $\sum_{j} z_j$, the sum of the jumps that occurred.

\paragraph{What characterises it.} Two ingredients, combined in its
\emph{intensity measure} $\nu(dz)\,dt$, the expected number of jumps of size
between $z$ and $z+dz$ during $dt$: an \emph{arrival rate}
$\lambda=\nu(\mathbb{R}^+)$, in jumps per unit time, and a \emph{size
distribution} $\nu(dz)/\lambda$. On the main panel, v1 estimates
$\lambda\approx5$ jumps per year, a mean size of $0.62$~nat and a maximum of
$3.21$~nats.

\paragraph{What $z$ is.} The size of an individual jump, in nats. Conceptually, it
is the mutual information of the news that caused it, in the sense of
Proposition~\ref{prop:conditioning}. Operationally, it is the change in $\JJ$
picked up by the four-sigma filter of Section~\ref{sec:program}. The two coincide
only approximately, and the paper does not claim to identify them.

\paragraph{Why the integral runs over $z>0$ only.} Because $\JJ$ measures an
entropy \emph{deficit} and, by Proposition~\ref{prop:conditioning}, information
can only increase that deficit in expectation. Within this framework nothing
creates a \emph{sudden increase} in entropy: entropy is restored only by the slow
relaxation of the first term, never by a jump. Hence one-sided jumps.

Two caveats, however. Proposition~\ref{prop:conditioning} is a statement
\emph{on average}: a particularly confusing realisation can raise entropy, so a
faithful specification would allow a small probability of negative jumps rather
than forbid them. And the Poisson measure assumes independent arrivals, whereas
the event study of v1 shows clear clustering of jumps; a self-exciting
Hawkes-type process would be more faithful.

Note finally that the diffusion coefficient is written $\sigma_\JJ$ rather than
$\eta$, the latter being reserved for the radius of stress scenarios
(Section~\ref{sec:stress}). Read on entropy itself rather than on $\JJ$, the sign
flips: entropy falls by jumps and recovers by diffusion. The observable
predictions are sharp, and they are what Section~\ref{sec:program} tests: the
distribution of $\Delta\JJ$ is right-skewed; large moves are up and small moves
down; and after a positive jump the index decays with half-life $\ln 2/\kappa$.

\subsection{Why fat tails, correlation and jumps arrive together}

The framework yields a prediction that no volatility or correlation model
produces on its own. It follows in four steps, without technical vocabulary.

\begin{enumerate}
\item Important news concerns \emph{all} assets at once: the same information
constrains each of them.
\item If the news is bad, all assets therefore fall \emph{on the same day}.
\item That day also appears, taken in isolation, in each asset's own history: for
that asset it is an unusually large loss.
\item Accumulating such days distorts each asset's return distribution: the loss
side lengthens and thickens more than the gain side.
\end{enumerate}

Step 2 is what the dependence channel measures: assets move together, measured
correlation rises, $-H^{dep}$ rises. Step 4 is what the shape channel measures:
extreme losses more numerous and larger than extreme gains, that is higher
kurtosis and negative skewness, so $\DD$ rises. The two channels rise together
because they describe the same day seen from two places: the cross-section for
one, each asset's distribution for the other.

In calm markets, notable moves are asset-specific: they do not align the
cross-section and so have no reason to lift both channels at the same time.

The testable form is a \emph{regime-dependent coupling}: the correlation between
$-\Delta H^{dep}$ and $\Delta\DD$ should be close to zero in calm markets and
strongly positive under stress. This is H6, and it is the paper's answer to the
question ``does the link between fat tails, correlation and entropy jumps
require an additional postulate?'': no. Given
Proposition~\ref{prop:conditioning}, the coupling follows from the common nature
of the news. Only the \emph{sign} requires P\ref{p:sign}.

\section{The H1--H8 empirical programme and what the nulls leave of it}
\label{sec:program}

This section summarises, without reproducing them, the estimators and tests of
arXiv~v1. It serves one purpose: to explain why, of all the theory in
Sections~\ref{sec:bridge} to~\ref{sec:dynamics}, only the dependence channel
makes it into the main paper.

\subsection{The eight hypotheses}

\begin{description}
\item[H1 (regime signature).] Under stress, $H^{dep}$ falls, $\DD$ rises, and so
$\JJ$ rises.
\item[H2 (partial compensation).] $\Delta H^{dep}_t=-\beta\,\Delta H^{vol}_t+u_t$
with $\beta>0$, and $\beta_{\text{stress}}>\beta_{\text{calm}}$.
\item[H3 (jump asymmetry).] Changes in $\JJ$ are one-sided, beyond what the same
estimator produces on data with no temporal structure.
\item[H4 (relaxation).] $\JJ$ mean-reverts with a half-life longer than the
estimator's own memory.
\item[H5 (skewness signature).] Positive jumps in $\JJ$ coincide with negative
skewness of the market factor (Postulate~P\ref{p:sign}).
\item[H6 (channel coupling).] Dependence and shape are weakly coupled in calm
markets and strongly coupled under stress.
\item[H7 (predictive power).] Today's entropy improves, out of sample, the
forecast of 21-day risk relative to volatility.
\item[H8 (composition of the crisis budget).] Grouped into episodes, jumps have
the dependence channel as their dominant contributor more often than a
memoryless resampling does.
\end{description}

\paragraph{Estimation, in one sentence.} Rolling 252-day windows,
\citet{ledoit2004} shrinkage covariance, marginal deficit via the
\citet{hyvarinen1998} approximation rather than the Edgeworth
expansion~\eqref{eq:edgeworth} (unusable on daily returns), stress regime
defined by a volatility threshold computed on the past only, jumps detected by a
four-sigma filter on $\Delta\JJ$.

\subsection{Why a calibrated null is indispensable}

The theoretical point that governs the whole reading of the results is the
following. $\JJ$ is a \emph{convex, non-negative} functional of second- to
fourth-order moments estimated over a window. When a fat-tailed observation
enters the window, the estimated moments jump, and the convex functional jumps
\emph{up} whatever the sign of the observation. As the observation ages
(exponential weighting) or leaves the window, the functional comes down
gradually. The measuring device therefore manufactures, on its own, the
sawtooth predicted by Proposition~\ref{prop:conditioning} and
Postulate~P\ref{p:relax}. An i.i.d.\ Gaussian simulation with no events at all
already yields an apparent half-life of 35~days and 37~positive jumps against
none negative under the four-sigma filter.

Hence two nulls, built by resampling the observed return matrix and running the
identical pipeline two hundred times: the \emph{i.i.d.} null (rows drawn
independently; marginal fat tails and cross-sectional dependence kept, all
chronology destroyed) and the 21-day \emph{block} null (which also keeps
volatility clustering). A statistic that clears both is evidence; one that
clears only one is not identified.

\subsection{The verdict}

\begin{table}[htbp]
\centering
\small
\begin{tabular}{@{}llrrrl@{}}
\toprule
& Statistic & Data & i.i.d. & block-21 & Verdict \\
\midrule
H1 & regime gap $H^{dep}$ & $-0.99$ & $-0.21$ & $-0.51$ & \textbf{clears both} \\
H1 & regime gap $\JJ$ & $1.93$ & $2.14$ & $1.97$ & fails both \\
H1 & regime gap $\DD$ & $0.27$ & $1.25$ & $0.65$ & wrong-signed \\
H2 & $\beta$ (stress) & $0.20$ & $0.05$ & $0.11$ & fails both \\
H3 & positive jumps & $158$ & $148.5$ & $133.2$ & block only \\
H3 & skewness of $\Delta\JJ$ & $12.26$ & $14.80$ & $12.97$ & wrong-signed \\
H4 & half-life (days) & $120.8$ & $92.1$ & $162.8$ & between the two \\
H4 & range of $\JJ$ & $18.95$ & $10.05$ & $13.77$ & \textbf{clears both} \\
H5 & $\mathrm{corr}(\Delta\JJ,\Delta S_{mkt})$ & $-0.24$ & $-0.04$ & $-0.07$ & fails both \\
H6 & coupling gap & $0.59$ & $0.77$ & $0.67$ & wrong-signed \\
H7 & out-of-sample gain & \multicolumn{3}{c}{none on two panels} & rejected \\
H8 & share of dependence-led episodes & $46\%$ & $19$--$33\%$ & n.i. & i.i.d.\ only \\
\bottomrule
\end{tabular}
\caption{Verdicts against the nulls (means over two hundred replications),
main panel of twenty S\&P~500 stocks, 1990--2022. Source:
\texttt{paper/tables/tab\_placebo\_sp500.tex} and the H7--H8 sections of v1.
n.i.: not identified, at every block length tested between 5 and 126~days.}
\label{tab:program}
\end{table}

Two statistics clear both nulls on the main panel, and only one replicates on
Kenneth French's 48 industry portfolios: the regime gap of the dependence
channel ($-3.65$~nats against $-0.56$ and $-2.46$). Phase~A then hardened it: it
holds in devolatilised (Forbes--Rigobon) and rank-based form, the latter being
the strongest against the block null ($p=0.03$ on the S\&P~500, $0.01$ on the
industries), whereas the \citet{kritzman2011} absorption ratio fails the block
null on both panels.

\subsection{What the verdict says about the theory}

One must separate what the verdict touches from what it does not.

\paragraph{What is untouched.} The statements of Sections~\ref{sec:bridge}
and~\ref{sec:channels} are identities and theorems. So is
Proposition~\ref{prop:conditioning}. No empirical result can refute them; it can
only show that the available estimators do not let us see them at work.

\paragraph{What is touched.} The two postulates. P\ref{p:relax} is supported only
by amplitude (the range of $\JJ$), not by a relaxation rate: the half-life lies
between the two nulls and its level depends on the null chosen. P\ref{p:sign} is
not established (H5). Channel coupling (H6), which Section~\ref{sec:dynamics}
derived from the common nature of news, is reproduced and even exceeded by
resampled data.

\paragraph{Why the dependence channel survives.} A plausible explanation,
consistent with Phase~A, lies in the form of the estimator.
$\tfrac12\log\det R=\tfrac12\sum_k\log\lambda_k$ weighs the whole spectrum of the
correlation matrix, including the small eigenvalues, which collapse when the
cross-section couples onto a few collective modes. A single large isolated return
barely moves those small eigenvalues. The shape channel $\DD$, by contrast,
reacts strongly to a single extreme observation, which is exactly the signature
the null reproduces. Concentration measures (absorption ratio, first-eigenvalue
share) look only at the top of the spectrum, where the block null catches up
with them.

\paragraph{Consequence for the main paper.} The indicator keeps the one
identified channel and drops $\DD$, and hence $\JJ$. The theory of
Sections~\ref{sec:channels} and~\ref{sec:dynamics} remains the justification of
\emph{why} this channel measures shared information (the Gaussian
multi-information, equation~\eqref{eq:ceiling} and Appendix~\ref{app:gauss}),
but the postulated dynamics is not used.

\section{The geometry of Kullback--Leibler-ball scenarios}
\label{sec:stress}

The main paper keeps the simplest use of this geometry: pricing a published
scenario in nats, which for a pure mean shift is half the squared Mahalanobis
distance. This section gives the full construction of which that use is a
special case, with the proofs that v1 only stated, and a new result on lifting a
portfolio scenario to full dimension.

\subsection{Divergence between two Gaussians}

\begin{proposition}[Closed form]
\label{prop:klgauss}
For $p_1=\mathcal N(\mu_1,\Sigma_1)$ and $p_0=\mathcal N(\mu_0,\Sigma_0)$ on
$\mathbb R^N$, with $\delta=\mu_1-\mu_0$,
\begin{equation}
\label{eq:klgauss}
\KL{p_1}{p_0}=\tfrac12\Big[\underbrace{\delta^\top\Sigma_0^{-1}\delta}_{\text{shift}}
+\underbrace{\operatorname{tr}(\Sigma_0^{-1}\Sigma_1)-N-\log\det(\Sigma_0^{-1}\Sigma_1)}_{\text{deformation}}\Big].
\end{equation}
\end{proposition}

\begin{proof}
$\KL{p_1}{p_0}=\mathbb E_{p_1}[\log p_1-\log p_0]$. Both logarithms are
quadratic; the trace trick of Appendix~\ref{app:gauss} gives
$\mathbb E_{p_1}[(X-\mu_1)^\top\Sigma_1^{-1}(X-\mu_1)]=N$ and
$\mathbb E_{p_1}[(X-\mu_0)^\top\Sigma_0^{-1}(X-\mu_0)]
=\operatorname{tr}(\Sigma_0^{-1}\Sigma_1)+\delta^\top\Sigma_0^{-1}\delta$.
Collect with the $\log\det$ terms.
\end{proof}

The shift term is half the squared Mahalanobis distance: this is the unit in
which the main paper audits published scenarios. The deformation term is
non-negative and zero if and only if $\Sigma_1=\Sigma_0$; writing $\nu_k$ for
the eigenvalues of $\Sigma_0^{-1}\Sigma_1$, it equals
$\tfrac12\sum_k(\nu_k-1-\log\nu_k)$, a sum of non-negative terms, one per
eigendirection.

\subsection{The projected ball and the exchange rate}

\begin{definition}[Entropy-ball scenario]
Given today's conditional law $p_0$ and a severity $\eta>0$, the severity-$\eta$
scenario for a risk functional $\rho$ is
\begin{equation}
\label{eq:ball}
p^\star_\eta=\arg\max_{p\,:\,\KL{p}{p_0}\le\eta}\rho(p).
\end{equation}
\end{definition}

For a linear portfolio $w^\top r$ with mean $m_0$ and standard deviation $s_0$
under $p_0$, set $u=(m-m_0)/s_0$ and $v=s/s_0$. By
Proposition~\ref{prop:klgauss} in one dimension,
\begin{equation}
\label{eq:g}
g(u,v)=\tfrac12\left(u^2+v^2-1-2\log v\right).
\end{equation}

\begin{proposition}[Solution of the projected ball]
\label{prop:ballsol}
For $\rho=-m+c\,s$ with $c\ge0$ (Gaussian VaR and expected shortfall are of this
form), the maximum of $\rho$ subject to $g(u,v)\le\eta$ is attained on the
boundary $g=\eta$ at a point satisfying
\begin{equation*}
u=-t,\qquad v-v^{-1}=c\,t,\qquad t>0 .
\end{equation*}
For $c=0$: $v=1$ and $u=-\sqrt{2\eta}$. \emph{A purely directional scenario of
$\eta$~nats is a move of $\sqrt{2\eta}$~sigmas}; one nat buys
$\sqrt2\approx1.41$~sigma, and $k$~sigmas cost $k^2/2$~nats.
\end{proposition}

\begin{proof}
$g$ is convex on $\mathbb R\times(0,\infty)$ (the sum of $u^2/2$ and
$(v^2-1-2\log v)/2$, whose second derivative is $1+v^{-2}>0$) and the objective
$s_0(-u+cv)$ is linear, so the Karush--Kuhn--Tucker conditions are sufficient.
The objective has no interior maximum, so the constraint binds. The Lagrangian
gives $-1=\lambda u$ and $c=\lambda(v-v^{-1})$ with $\lambda>0$; setting
$t=1/\lambda$ yields the stated system, and $t$ is pinned down by
$g(-t,v(t))=\eta$. For $c=0$, $v-v^{-1}=0$ forces $v=1$ and $g=u^2/2=\eta$.
\end{proof}

\subsection{From the portfolio to full dimension}

The ball~\eqref{eq:ball} is solved on the portfolio's law, but a stress test must
say how each asset moves. Two facts settle the question.

\begin{proposition}[Data processing]
\label{prop:dpi}
For any law $p$ on $\mathbb R^N$, with $p^w$ and $p_0^w$ the laws of $w^\top r$
under $p$ and $p_0$,
$\KL{p^w}{p_0^w}\le\KL{p}{p_0}$.
\end{proposition}

This is the data-processing inequality \citep{cover2006} applied to the map
$r\mapsto w^\top r$. The \emph{ambient} price of a scenario (in full dimension)
is therefore always at least its \emph{portfolio} price (in projection). The
classical-bundle table of v1 bears this out line by line ($9.97$ against
$7.24$~nats for the standard bundle).

\begin{proposition}[Minimal lift]
\label{prop:lift}
Let $p_0=\mathcal N(\mu_0,\Sigma_0)$, $m_0=w^\top\mu_0$, $s_0^2=w^\top\Sigma_0w$,
and a target $(m,s)$ with $s>0$. Among Gaussians $\mathcal N(\mu_1,\Sigma_1)$
with $w^\top\mu_1=m$ and $w^\top\Sigma_1w=s^2$, the divergence $\KL{p_1}{p_0}$
is minimised by
\begin{equation}
\label{eq:lift}
\mu_1=\mu_0+u\,s_0\,\beta,\qquad
\Sigma_1=\Sigma_0+(v^2-1)\,s_0^2\,\beta\beta^\top,
\qquad\beta=\frac{\Sigma_0w}{s_0^2},
\end{equation}
and the minimum equals exactly $g(u,v)$, the portfolio price.
\end{proposition}

\begin{proof}
By Proposition~\ref{prop:dpi}, every candidate costs at least the price of its
projection, which is $g(u,v)$ since the projection of a Gaussian is Gaussian
with mean $m$ and standard deviation $s$. It therefore suffices to check that
\eqref{eq:lift} attains this bound. We have $w^\top\beta=1$, so
$w^\top\mu_1=m_0+us_0=m$ and $w^\top\Sigma_1w=s_0^2+(v^2-1)s_0^2=s^2$. In
\eqref{eq:klgauss}: $\delta^\top\Sigma_0^{-1}\delta=u^2s_0^2\,\beta^\top\Sigma_0^{-1}\beta=u^2$;
$\operatorname{tr}(\Sigma_0^{-1}\Sigma_1)=N+(v^2-1)s_0^2\,\beta^\top\Sigma_0^{-1}\beta=N+v^2-1$;
and by the matrix determinant lemma,
$\det(\Sigma_0^{-1}\Sigma_1)=1+(v^2-1)s_0^2\,\beta^\top\Sigma_0^{-1}\beta=v^2$.
Hence $\KL{p_1}{p_0}=\tfrac12(u^2+v^2-1-2\log v)=g(u,v)$. The matrix $\Sigma_1$ is
positive definite for every $v>0$.
\end{proof}

\paragraph{Asset-by-asset reading.} The vector $\beta$ collects the assets'
betas on the portfolio, $\beta_i=\operatorname{Cov}(r_i,w^\top r)/s_0^2$. Write
$r=\beta f+e$ with $f=w^\top r$ and $e$ uncorrelated with $f$ under $p_0$. The
scenario~\eqref{eq:lift} multiplies the standard deviation of the portfolio
factor $f$ by $v$, shifts its mean by $u$~sigmas, and leaves the law of $e$
unchanged. Each asset's variance rises by $(v^2-1)s_0^2\beta_i^2$ and its mean
moves by $us_0\beta_i$: high-beta assets take the shock, and the correlations
among them rise. It is a one-factor scenario, and it is the least informative of
all scenarios that reach the target severity.

\begin{remark}[Correction to v1]
arXiv~v1 identified this minimum through its first-order conditions,
$\Sigma_1^{-1}=\Sigma_0^{-1}-\lambda ww^\top$, and judged it ``not interpretable
asset by asset''. By the Sherman--Morrison formula, this rank-one perturbation of
the precision matrix \emph{is}~\eqref{eq:lift}, a rank-one perturbation of the
covariance along $\Sigma_0w$, and it reads as above. Since v1 never uses it, no
number depends on it.
\end{remark}

\begin{corollary}[Cost of uniform dilation]
\label{cor:uniform}
Uniform dilation $\Sigma_1=v^2\Sigma_0$, the v1 default, has a deformation
component equal to $N\cdot\tfrac12(v^2-1-2\log v)$, that is $N$~times that of the
minimal lift.
\end{corollary}

\begin{proof}
Every eigenvalue of $\Sigma_0^{-1}\Sigma_1$ equals $v^2$; apply the spectral form
of the deformation term.
\end{proof}

For the ten-year scenario on the main panel ($v=1.63$, $N=20$), the variance
component is $0.34$~nat at the minimum and $20\times0.34\approx6.8$ under uniform
dilation, which is the figure in the v1 composition table. The corollary
explains the v1 finding that the compositions calibrated on H1 ($2.88$~nats) and
refined by H8 ($0.71$) cost \emph{less} than the default: the default deforms
every eigendirection equally, whereas a composition that concentrates the shock
on the common mode moves towards the one-factor scenario~\eqref{eq:lift}, whose
total cost is $0.41$~nat including the mean shift. That uniform dilation costs
$N$~times the minimum in variance is a general result; that the calibrated
compositions fall between the two remains a finding on these panels, but it
becomes expected rather than surprising.

\subsection{Composing a given severity}

\begin{proposition}[One severity, many compositions]
\label{prop:composition}
If $w^\top\Sigma_aw=w^\top\Sigma_bw=s^2$, then for every $\alpha\in[0,1]$,
$\Sigma_\alpha=(1-\alpha)\Sigma_a+\alpha\Sigma_b$ satisfies
$w^\top\Sigma_\alpha w=s^2$.
\end{proposition}

\begin{proof}
$w^\top\Sigma w$ is linear in $\Sigma$.
\end{proof}

Portfolio severity therefore does not constrain the composition. The ambient
price, for its part, is convex in $\Sigma_1$ (the $-\log\det$ term is convex and
the rest is linear), so along a mixture it is bounded by the mixture of prices,
which guarantees that an intermediate composition never costs more than the
dearer of its two endpoints.

\paragraph{A physical ceiling on correlation.} With marginal volatilities fixed,
$w^\top\Sigma w$ is maximised over valid correlation matrices at
$\rho_{ij}=1$, where it equals $(\sum_iw_i\sigma_i)^2$. On the main panel this
ceiling is only about $1.5$~times today's volatility: beyond the five-year rung,
no rise in correlation alone suffices, and some rise in volatilities is
unavoidable.

\subsection{Calibrating the radius, and the limit of that calibration}

The radius is read off the information flow of the projection,
$I^{(h)}_t=\KL{p^w_t}{p^w_{t-h}}$, whose quantiles (on non-overlapping 21-day
blocks) convert a return period into a radius. Reading the radius off the
$N$-dimensional flow overstates it by an order of magnitude, by
Proposition~\ref{prop:dpi}. Calibrated this way, the v1 ladder gives a ten-year
volatility multiplier of $1.63$, while the ten-year quantile of the realised
one-month volatility rise is $3.2$. The ladder is a coherent, reproducible rule;
its return periods are not validated. This is one of the reasons the main paper
keeps the audit (pricing a given scenario) and drops generation (building a
scenario with a given return period).

\subsection{Reverse stress testing}

Asking for the cheapest scenario that produces a given loss $L$ inverts the ball.
In the projected Gaussian setting, Proposition~\ref{prop:ballsol} reads
backwards. In the non-parametric setting, the solution is the exponential
tilting of the sample \citep{csiszar1975}, whose effective sample size collapses
as the radius grows: at $9$~nats it is $1.0$ on the main panel, meaning that the
tilted law is nothing but the worst observed day. Any implementation should
report this effective size and switch to the Gaussian ball as soon as it drops
below a few dozen.

\section{The bridge to the main paper: block decomposition}
\label{sec:blocks}

This section connects the theory above to the indicator of the main paper. It
establishes the identity behind the ``contagion map'' and shows where the
dependence channel sits within the structural index.

\subsection{The identity}

Partition the $N$ series into $B$ blocks (for instance the CISS segments: money
market, bonds, equities, financial intermediaries, foreign exchange), and write
$R_{bb}$ for the diagonal blocks of the correlation matrix $R$ and
$R_D=\operatorname{diag}(R_{11},\dots,R_{BB})$.

\begin{proposition}[Within/between-block decomposition]
\label{prop:blocks}
\begin{equation}
\label{eq:blocks}
\underbrace{-\tfrac12\log\det R}_{\text{total dependence}}
\;=\;\sum_{b=1}^{B}\underbrace{\left(-\tfrac12\log\det R_{bb}\right)}_{\text{within block }b}
\;+\;\underbrace{\left(-\tfrac12\log\det\big(R_D^{-1/2}RR_D^{-1/2}\big)\right)}_{\text{between blocks}} ,
\end{equation}
and each of the $B+1$ terms is non-negative. The between-block term is the
Gaussian multi-information among blocks,
$\KL{\mathcal N(0,R)}{\bigotimes_b\mathcal N(0,R_{bb})}$; it is zero if and only
if the blocks are pairwise uncorrelated.
\end{proposition}

\begin{proof}
$R=R_D^{1/2}\,\big(R_D^{-1/2}RR_D^{-1/2}\big)\,R_D^{1/2}$ and the determinant is
multiplicative, so $\log\det R=\log\det R_D+\log\det(R_D^{-1/2}RR_D^{-1/2})$,
with $\log\det R_D=\sum_b\log\det R_{bb}$. Each $R_{bb}$ is a correlation matrix,
so $-\tfrac12\log\det R_{bb}\ge0$ by Hadamard's inequality
(Appendix~\ref{app:gauss}). The matrix $\tilde R=R_D^{-1/2}RR_D^{-1/2}$ has
identity diagonal blocks, hence a unit diagonal, and Hadamard applies again.
Finally Proposition~\ref{prop:klgauss} with $\mu_1=\mu_0$, $\Sigma_1=R$,
$\Sigma_0=R_D$ gives
$\tfrac12[\operatorname{tr}(R_D^{-1}R)-N-\log\det(R_D^{-1}R)]$, and
$\operatorname{tr}(R_D^{-1}R)=\sum_b\operatorname{tr}(I)=N$.
\end{proof}

\paragraph{Two blocks: canonical correlations.} For $B=2$, writing
$C=R_{11}^{-1/2}R_{12}R_{22}^{-1/2}$ and $\rho_1,\dots,\rho_k$ for its singular
values (the canonical correlations between the two blocks),
$\det\tilde R=\det(I-CC^\top)=\prod_j(1-\rho_j^2)$ and the between-block term
equals $-\tfrac12\sum_j\log(1-\rho_j^2)$. This is the Gaussian mutual
information between the two segments, the sum of the information carried by each
canonical pair. For $B>2$, identity~\eqref{eq:blocks} can be applied in nested
fashion and lets one attribute contagion to pairs or groups of segments by the
chain rule.

\subsection{Where the dependence channel sits in the structural index}

The chain rule of Appendix~\ref{app:index} writes $\JJ=\sum_i\Dun(x_i)+I(x)$.
Replacing the multi-information $I(x)$ by its Gaussian part $-H^{dep}$ and
applying~\eqref{eq:blocks}:
\begin{equation}
\label{eq:jblocks}
\JJ\;\approx\;\underbrace{\sum_i\Dun(x_i)}_{\text{shape}}
\;+\;\underbrace{\sum_b\left(-\tfrac12\log\det R_{bb}\right)}_{\text{within}}
\;+\;\underbrace{\left(-\tfrac12\log\det\tilde R\right)}_{\text{between}} .
\end{equation}
The main paper keeps the last two terms and drops the first, which is the one
the null reproduces (Section~\ref{sec:program}). The entropic CISS is thus the
structural index stripped of its shape channel, and its decomposition is exact,
with no approximation, because it concerns only the Gaussian part.

\subsection{Three properties useful for the indicator}

\paragraph{Scale invariance.} The terms of~\eqref{eq:blocks} depend only on
$R$: they are invariant to any series-by-series rescaling, which allows markets
with very different volatilities to be aggregated without ad hoc weighting. The
\citet{forbes2002} devolatilisation therefore acts not on the formula but on the
bias of the \emph{estimator} of $R$ under heteroskedasticity.

\paragraph{Rank version.} Replacing $R$ by the correlation of normal scores
$\Phi^{-1}(\hat F_i(x_i))$ estimates the multi-information of a Gaussian copula.
The quantity becomes invariant to any monotone transformation of the marginals,
which explains its robustness to tails observed in Phase~A.

\paragraph{Normalisation in $N$.} For a constant correlation $\rho$,
$\det R=(1-\rho)^{N-1}\big(1+(N-1)\rho\big)$, so $-H^{dep}$ grows roughly linearly
in $N$. Dividing by $N-1$, or reporting the equivalent constant correlation
$\rho^\star$ that gives the same determinant, makes the indicator comparable
across panels. For small $\rho^\star$,
$-H^{dep}\approx\tfrac14N(N-1)\rho^{\star2}$.

\section{Monitoring indicators and insurance solvency}
\label{sec:monitoring}

\paragraph{Two indicators from v1.} The jump/diffusion split delivered two daily
quantities: the standardised positive-jump component of $\JJ$ (how much
information arrived today) and $\JJ_t$ relative to its rolling maximum (how much
has not yet been absorbed). Neither improves the volatility forecast (H7), and
both rest on statistics the null does not identify. They are not carried over.

\paragraph{Insurance.} Three links remain valid in principle. An entropy-ball
scenario produces a complete stressed law, a natural object for the ORSA. The
decomposition of classical bundles bears on the Solvency~II standard formula,
whose diversification credit rests on a fixed correlation matrix. And entropic
risk measures are coherent \citep{ahmadijavid2012,artzner1999}. Nothing has been
tested on liabilities, however, and the short, low-frequency panels typical of
insurance data would strain every estimator used here.

\section{Open questions}
\label{sec:open}

\begin{enumerate}
\item \textbf{Self-exciting arrivals.} The jump clustering seen in the v1 event
study contradicts the Poisson measure in~\eqref{eq:ou}. A Hawkes process is the
next specification.
\item \textbf{Positivity.} $\JJ\ge0$ whereas an Ornstein--Uhlenbeck process with
a Brownian term can turn negative. One should either drop $\sigma_\JJ$ (the
pure form of \citealt{bns2001}) or model $\log\JJ$.
\item \textbf{Non-linear tail dependence.} $\JJ$ and the dependence channel see
only the Gaussian part of the multi-information. Estimating the full copula
multi-information (for example by nearest neighbours on ranks) would measure
what is missing, at the price of a much higher estimation variance.
\item \textbf{Tail of the information flow.} Beyond about twenty years, the
quantiles of $I^{(h)}$ are extrapolated. A generalised Pareto fit would be the
disciplined correction.
\item \textbf{Minimal lift and published scenarios.}
Proposition~\ref{prop:lift} provides a yardstick: for any published scenario,
the ratio of its ambient price to the minimal price $g(u,v)$ of its projection
measures how much information it spends outside the direction that matters to
the portfolio. This extends the audit of the main paper directly.
\end{enumerate}

\appendix

\section{Units, invariance, and what can be computed}
\label{app:units}

\paragraph{The nat.} Entropy $H=-\sum_i p_i\log p_i$ is defined only once the
base of the logarithm is chosen, and that choice \emph{is} the choice of unit.
With $\log_2$ one gets bits (shannons), with $\log_{10}$ dits, and with the
natural logarithm \emph{nats}. One nat is the information delivered by an
observation that divides the space of possibilities by $e\approx2.718$, where a
bit divides it by two. The conversion is $1~\text{nat}=1/\log 2\approx1.4427$~bit.
This document uses the natural logarithm throughout: entropies, deficits,
divergences, scenario radii and the structural index are in nats.

Three operational readings, from the most useful to the most abstract. First the
exchange rate established in Section~\ref{sec:stress}: a purely directional
scenario of $\eta$ nats is a move of $\sqrt{2\eta}$ sigmas, so $k$~sigmas cost
$k^2/2$ nats. Next, coding: $\KL{p}{q}$ nats is the expected excess code length
when data with law $p$ are encoded with a code optimal for $q$. Finally,
hypothesis testing: a divergence of one nat per observation means that the
likelihood ratio in favour of $p$ grows on average by a factor $e$ with each
observation.

\paragraph{$h$ depends on units, divergences do not.} The differential entropy
of a continuous variable satisfies, for any invertible affine transformation,
\begin{equation}
\label{eq:h-scaling}
h(aX+b) \;=\; h(X) + \log|a| ,
\qquad\text{and in dimension } N,\quad h(AX+b)=h(X)+\log|\det A| .
\end{equation}
Expressing returns in percent rather than decimals therefore adds $N\log 100$ to
$h$: the \emph{level} of $h$, and hence that of $H^{vol}$ and $H^{cov}$, has no
interpretable absolute value. By contrast $\KL{p}{q}$, $I(X;Y)$, $\DD$ and $\JJ$
are invariant under any invertible change of variable, the Jacobians cancelling
between numerator and denominator of the logarithm. \emph{A figure in nats is
comparable across datasets only if it is a divergence.} This is the structural
reason the paper builds a scale-free index.

\paragraph{How is $h(X)$ computed?} In the Gaussian case, in closed form
(Appendix~\ref{app:gauss}). In general, it cannot be without knowing $p$.
Non-parametric estimators exist (\citet{vasicek1976} spacings,
Kozachenko--Leonenko nearest neighbours, kernel methods) but all degrade quickly
beyond three or four dimensions. Estimating a differential entropy in
dimension~20 on a 252-day window is out of reach.

The paper therefore \emph{never} computes $h(p)$. It computes $h(q)$ in closed
form from $\widehat\Sigma$, and the deficit $\DD$ through an approximation of
the \emph{difference} of marginal entropies. A difference of entropies is
dimensionless and estimable; a bare differential entropy is not. The whole
architecture (decompose rather than estimate directly) follows from this
constraint.

\section{Gaussian entropy: proof of equation~\eqref{eq:ceiling}}
\label{app:gauss}

\begin{proof}[Step A: entropy of a multivariate Gaussian]
For $q=\mathcal N(\mu,\Sigma)$ on $\mathbb{R}^N$,
\[
\log q(x) = -\tfrac12\left[N\log 2\pi + \log\det\Sigma + (x-\mu)'\Sigma^{-1}(x-\mu)\right],
\]
so
\[
h(q) = -\mathbb{E}_q[\log q] = \tfrac12 N\log 2\pi + \tfrac12\log\det\Sigma
+ \tfrac12\,\mathbb{E}\!\left[(X-\mu)'\Sigma^{-1}(X-\mu)\right].
\]
The last expectation follows from the trace trick: a scalar equals its trace,
and the trace is cyclic, hence
\[
\mathbb{E}\!\left[(X-\mu)'\Sigma^{-1}(X-\mu)\right]
= \mathbb{E}\!\left[\operatorname{tr}\!\left(\Sigma^{-1}(X-\mu)(X-\mu)'\right)\right]
= \operatorname{tr}(\Sigma^{-1}\Sigma) = \operatorname{tr}(I_N) = N .
\]
Collecting terms, $h(q)=\tfrac12\left[N\log2\pi+\log\det\Sigma+N\right]
=\tfrac12\log\!\left((2\pi e)^N\det\Sigma\right)$.
\end{proof}

\begin{proof}[Step B: factorising the determinant]
With $\Sigma=DRD$, $D=\operatorname{diag}(\sigma_1,\dots,\sigma_N)$ and $R$ the
correlation matrix, multiplicativity of the determinant gives
$\det\Sigma=\det(D)\det(R)\det(D)=\left(\prod_i\sigma_i\right)^2\det R$.
\end{proof}

\begin{proof}[Step C: substitution]
$\tfrac12\log\det\Sigma
=\tfrac12\log\!\left[\left(\prod_i\sigma_i\right)^2\det R\right]
=\sum_i\log\sigma_i+\tfrac12\log\det R$, hence
$h(q)=\tfrac N2\log(2\pi e)+\sum_i\log\sigma_i+\tfrac12\log\det R$, which is
equation~\eqref{eq:ceiling}.
\end{proof}

\paragraph{Sign of the dependence channel.} Hadamard's inequality gives
$\det R\le\prod_i R_{ii}=1$ for a positive semi-definite matrix. A more telling
proof goes through the spectrum: $\det R=\prod_k\lambda_k$ with
$\sum_k\lambda_k=\operatorname{tr}R=N$, and the arithmetic--geometric mean
inequality gives $\prod_k\lambda_k\le\left(\sum_k\lambda_k/N\right)^N=1$, with
equality if and only if every eigenvalue equals~1, that is $R=I$. Hence
$H^{dep}\le0$: \emph{correlation always destroys entropy}.

In the bivariate case, $\det R=1-\rho^2$ and $H^{dep}=\tfrac12\log(1-\rho^2)$;
for $\rho=0.9$ this is $-0.83$~nat. Note that $-H^{dep}$ is then exactly the
mutual information between the two assets: correlation is shared information,
hence destroyed entropy.

\paragraph{Reading note on the tables.} The $H^{cov}$ columns of the v1 results
tables report $H^{vol}+H^{dep}$ and \textbf{omit the constant}
$\tfrac N2\log(2\pi e)$, which equals $28.38$~nats for $N=20$. This constant
offset is harmless since only differences between regimes and changes over time
are interpreted, but it explains why the reported levels are negative.

\section{The entropy deficit: why $\DD=h(q)-h(p)$}
\label{app:deficit}

The Kullback--Leibler divergence between a law $p$ and a reference law $q$ is
$\KL{p}{q}=\int p\log(p/q)$, non-negative and zero if and only if $p=q$ almost
everywhere. It is neither symmetric nor equipped with a triangle inequality: it
is not a distance, hence the word divergence.

\begin{proposition}
\label{prop:deficit}
If $q=\mathcal N(\mu,\Sigma)$ has \emph{the first two moments of $p$}, then
$\KL{p}{q}=h(q)-h(p)$.
\end{proposition}

\begin{proof}
$\KL{p}{q}=-h(p)-\mathbb{E}_p[\log q(X)]$. Now $\log q$ is exactly quadratic in
$x$, so its expectation depends only on the mean and covariance of the law under
which it is evaluated. Since $p$ and $q$ share these two moments,
$\mathbb{E}_p\!\left[(X-\mu)'\Sigma^{-1}(X-\mu)\right]=\operatorname{tr}(\Sigma^{-1}\Sigma)=N
=\mathbb{E}_q[\cdot]$, hence $\mathbb{E}_p[\log q]=\mathbb{E}_q[\log q]=-h(q)$ and
so $\KL{p}{q}=h(q)-h(p)$.
\end{proof}

The moment-matching assumption is essential: without it the identity fails. This
is the whole content of Proposition~\ref{prop:maxent}, read as a measurement
rather than an existence theorem. Economically, $\DD$ is the price in nats of the
lie ``the market is Gaussian'': the apparent uncertainty that the Gaussian
description overstates, because the true law is more concentrated than its
covariance alone suggests.

\section{The structural index is a divergence}
\label{app:index}

\begin{proposition}
\label{prop:index}
With $q_i=\mathcal N(\mu_i,\sigma_i^2)$ the Gaussian marginals with the same
scales as $p$,
\[
\KL{p}{\textstyle\prod_i q_i} \;=\; \DD - H^{dep} \;=\; \JJ .
\]
\end{proposition}

\begin{proof}
By definition,
\[
\begin{aligned}
\KL{p}{\textstyle\prod_i q_i} &= -h(p) - \mathbb{E}_p\!\left[\textstyle\sum_i\log q_i(X_i)\right]
= -h(p) + \sum_i\left[\tfrac12\log(2\pi\sigma_i^2)+\tfrac12\right]\\
&= -h(p) + \sum_i\log\sigma_i + \tfrac N2\log(2\pi e),
\end{aligned}
\]
using $\mathbb{E}_p[(X_i-\mu_i)^2]=\sigma_i^2$, true by construction. Substituting
equation~\eqref{eq:three}, namely
$h(p)=\tfrac N2\log(2\pi e)+\sum_i\log\sigma_i+H^{dep}-\DD$, the scale terms and
the constant cancel and $\DD-H^{dep}$ remains.
\end{proof}

This is an exact identity, not an approximation, and it is what justifies
Definition~\ref{def:J}. Removing the scale channel from \eqref{eq:three} might
look like a fudge; the object that remains turns out to be a proper
Kullback--Leibler divergence, and so inherits three properties for free: it is
non-negative; it is zero if and only if returns are independent \emph{and}
Gaussian, which provides a well-defined rather than arbitrary reference point;
and it is invariant to asset-by-asset changes of units, since $p$ and
$\prod_i q_i$ transform in the same way.

It is useful to set this identity beside the general chain rule
$\KL{p}{\prod_i q_i}=\sum_i\KL{p_i}{q_i}+I(x)$, where
$I(x)=\KL{p}{\prod_i p_i}$ is the multi-information. The first term is the sum of
marginal deficits $\sum_i\Dun(x_i)$, the second the total dependence. The
estimator used in v1 estimates the first term with the \citet{hyvarinen1998}
approximation and replaces the second with its linear part $-H^{dep}$;
non-linear tail dependence is therefore omitted from $\JJ$ and measured
separately through a co-exceedance rate.

\section{Mutual information: the two propositions and an example}
\label{app:mi}

\paragraph{Proof of Proposition~\ref{prop:conditioning}.} By definition of
conditional entropy, $\mathbb{E}_Y[h(X\mid Y)]=h(X,Y)-h(Y)$, and by
\eqref{eq:mi}, $I(X;Y)=h(X)+h(Y)-h(X,Y)$. Adding the two,
$\mathbb{E}_Y[h(X\mid Y)]=h(X)-I(X;Y)$. Non-negativity of $I$ follows from that of
the Kullback--Leibler divergence, of which it is a special case, with equality if
and only if $p_{XY}=p_X\otimes p_Y$.

\paragraph{Proof of Proposition~\ref{prop:constraint}.} Immediate: an extra
constraint can only shrink the feasible set, and the supremum of a functional
over a smaller set cannot increase.

\paragraph{Why the two are not equivalent.} The first quantifies but holds only
on average; the second holds realisation by realisation but concerns the ceiling
and quantifies nothing. A discrete counterexample suffices to show that a
particular realisation can \emph{raise} entropy. Let $X$ be binary with
$\mathbb{P}(X=a)=0.99$, so $H(X)=0.056$~nat. Let a signal $Y$, with probability
$0.98$, remove all doubt ($H(X\mid Y=0)=0$), and with probability $0.02$ lead to
$\mathbb{P}(X=a\mid Y=1)=1/2$, so $H(X\mid Y=1)=\log 2=0.693$~nat, twelve times the
prior entropy. The average is nonetheless
$0.98\times0+0.02\times0.693=0.014$~nat, well below $0.056$, and
$I(X;Y)=0.042$~nat. This is exactly why the paper predicts a jump component that
is \emph{mainly}, not exclusively, one-sided.

\paragraph{Worked Gaussian example.} Let $X\sim\mathcal N(0,\sigma^2)$ be
tomorrow's return and $Y=X+\varepsilon$ a noisy signal,
$\varepsilon\sim\mathcal N(0,\sigma_\varepsilon^2)$ independent. Then
\begin{equation}
\label{eq:mi-gauss}
I(X;Y)=\tfrac12\log\!\left(1+\frac{\sigma^2}{\sigma_\varepsilon^2}\right),
\qquad
\operatorname{Var}(X\mid Y)=\frac{\sigma^2\sigma_\varepsilon^2}{\sigma^2+\sigma_\varepsilon^2},
\end{equation}
the posterior variance not depending on the realisation $y$. With
$\sigma=1\%$ a day:

\begin{center}
\small
\begin{tabular}{@{}lrrrr@{}}
\toprule
Signal quality & $\sigma_\varepsilon$ & $I(X;Y)$ (nats) & posterior vol. & $\sqrt{2\eta}$ equivalent \\
\midrule
sharp    & $0.5\%$ & $0.805$ & $0.45\%$ & $1.27\sigma$ \\
moderate & $1.0\%$ & $0.347$ & $0.71\%$ & $0.83\sigma$ \\
vague    & $3.0\%$ & $0.053$ & $0.95\%$ & $0.32\sigma$ \\
\bottomrule
\end{tabular}
\end{center}

A factor of fifteen separates the size of the entropy jump caused by decisive
news from that caused by a rumour. The last column translates each jump onto the
severity scale of Section~\ref{sec:stress}: it is the directional move a
scenario of the same radius would produce.

\begin{thebibliography}{99}
\small

\bibitem[Ahmadi-Javid(2012)]{ahmadijavid2012} Ahmadi-Javid, A. (2012). Entropic Value-at-Risk: A New Coherent Risk Measure. \textit{Journal of Optimization Theory and Applications}, 155(3), 1105--1123.

\bibitem[Almog \& Shmueli(2019)]{almog2019} Almog, A., \& Shmueli, E. (2019). Structural Entropy: Monitoring Correlation-Based Networks Over Time With Application To Financial Markets. \textit{Scientific Reports}, 9, 10832.

\bibitem[Ané \& Geman(2000)]{ane2000} Ané, T., \& Geman, H. (2000). Order Flow, Transaction Clock, and Normality of Asset Returns. \textit{Journal of Finance}, 55(5), 2259--2284.

\bibitem[Artzner et al.(1999)]{artzner1999} Artzner, P., Delbaen, F., Eber, J.-M., \& Heath, D. (1999). Coherent Measures of Risk. \textit{Mathematical Finance}, 9(3), 203--228.

\bibitem[Barndorff-Nielsen \& Shephard(2001)]{bns2001} Barndorff-Nielsen, O. E., \& Shephard, N. (2001). Non-Gaussian Ornstein--Uhlenbeck-based Models and Some of Their Uses in Financial Economics. \textit{Journal of the Royal Statistical Society B}, 63(2), 167--241.

\bibitem[Barndorff-Nielsen \& Shephard(2004)]{bns2004} Barndorff-Nielsen, O. E., \& Shephard, N. (2004). Power and Bipower Variation with Stochastic Volatility and Jumps. \textit{Journal of Financial Econometrics}, 2(1), 1--37.

\bibitem[Black(1976)]{black1976} Black, F. (1976). Studies of Stock Price Volatility Changes. \textit{Proceedings of the Business and Economic Statistics Section, American Statistical Association}, 177--181.

\bibitem[Christie(1982)]{christie1982} Christie, A. A. (1982). The Stochastic Behavior of Common Stock Variances: Value, Leverage and Interest Rate Effects. \textit{Journal of Financial Economics}, 10(4), 407--432.

\bibitem[Doob(1942)]{doob1942} Doob, J. L. (1942). The Brownian Movement and Stochastic Equations. \textit{Annals of Mathematics}, 43(2), 351--369.

\bibitem[Engle \& Ng(1993)]{engle1993} Engle, R. F., \& Ng, V. K. (1993). Measuring and Testing the Impact of News on Volatility. \textit{Journal of Finance}, 48(5), 1749--1778.

\bibitem[Nelson(1991)]{nelson1991} Nelson, D. B. (1991). Conditional Heteroskedasticity in Asset Returns: A New Approach. \textit{Econometrica}, 59(2), 347--370.

\bibitem[Bollerslev \& Ghysels(1996)]{bollerslev1996} Bollerslev, T., \& Ghysels, E. (1996). Periodic Autoregressive Conditional Heteroscedasticity. \textit{Journal of Business \& Economic Statistics}, 14(2), 139--151.

\bibitem[Breuer \& Csiszár(2013)]{breuer2013} Breuer, T., \& Csiszár, I. (2013). Systematic Stress Tests with Entropic Plausibility Constraints. \textit{Journal of Banking \& Finance}, 37(5), 1552--1559.

\bibitem[Caraiani(2018)]{caraiani2018} Caraiani, P. (2018). Modeling the Comovement of Entropy between Financial Markets. \textit{Entropy}, 20(6), 417.

\bibitem[Clark(1973)]{clark1973} Clark, P. K. (1973). A Subordinated Stochastic Process Model with Finite Variance for Speculative Prices. \textit{Econometrica}, 41(1), 135--155.

\bibitem[Cover \& Thomas(2006)]{cover2006} Cover, T. M., \& Thomas, J. A. (2006). \textit{Elements of Information Theory}, 2nd ed. Wiley.

\bibitem[Csiszár(1975)]{csiszar1975} Csiszár, I. (1975). $I$-Divergence Geometry of Probability Distributions and Minimization Problems. \textit{Annals of Probability}, 3(1), 146--158.

\bibitem[Diebold \& Mariano(1995)]{dm1995} Diebold, F. X., \& Mariano, R. S. (1995). Comparing Predictive Accuracy. \textit{Journal of Business \& Economic Statistics}, 13(3), 253--263.

\bibitem[Engle(2002)]{engle2002} Engle, R. (2002). Dynamic Conditional Correlation. \textit{Journal of Business \& Economic Statistics}, 20(3), 339--350.

\bibitem[Forbes \& Rigobon(2002)]{forbes2002} Forbes, K. J., \& Rigobon, R. (2002). No Contagion, Only Interdependence: Measuring Stock Market Comovements. \textit{Journal of Finance}, 57(5), 2223--2261.

\bibitem[Holló, Kremer \& Lo Duca(2012)]{hollo2012} Holló, D., Kremer, M., \& Lo Duca, M. (2012). CISS -- A Composite Indicator of Systemic Stress in the Financial System. \textit{ECB Working Paper} No.~1426.

\bibitem[Kritzman et al.(2011)]{kritzman2011} Kritzman, M., Li, Y., Page, S., \& Rigobon, R. (2011). Principal Components as a Measure of Systemic Risk. \textit{Journal of Portfolio Management}, 37(4), 112--126.

\bibitem[Glasserman \& Xu(2014)]{glasserman2014} Glasserman, P., \& Xu, X. (2014). Robust Risk Measurement and Model Risk. \textit{Quantitative Finance}, 14(1), 29--58.

\bibitem[Hyvärinen(1998)]{hyvarinen1998} Hyvärinen, A. (1998). New Approximations of Differential Entropy for Independent Component Analysis and Projection Pursuit. \textit{Advances in Neural Information Processing Systems}, 10, 273--279. (Constants as in Hyvärinen \& Oja, 2000, \textit{Neural Networks} 13, 411--430.)

\bibitem[Jaynes(1957)]{jaynes1957} Jaynes, E. T. (1957). Information Theory and Statistical Mechanics. \textit{Physical Review}, 106, 620--630.

\bibitem[Kyle(1985)]{kyle1985} Kyle, A. S. (1985). Continuous Auctions and Insider Trading. \textit{Econometrica}, 53(6), 1315--1335.

\bibitem[Ledoit \& Wolf(2004)]{ledoit2004} Ledoit, O., \& Wolf, M. (2004). A Well-Conditioned Estimator for Large-Dimensional Covariance Matrices. \textit{Journal of Multivariate Analysis}, 88(2), 365--411.

\bibitem[Mandelbrot \& Taylor(1967)]{mandelbrot1967} Mandelbrot, B., \& Taylor, H. M. (1967). On the Distribution of Stock Price Differences. \textit{Operations Research}, 15(6), 1057--1062.

\bibitem[Monroe(1978)]{monroe1978} Monroe, I. (1978). Processes That Can Be Embedded in Brownian Motion. \textit{Annals of Probability}, 6(1), 42--56.

\bibitem[Papla \& Siedlecki(2024)]{papla2024} Papla, D., \& Siedlecki, R. (2024). Entropy as a Tool for the Analysis of Stock Market Efficiency During Periods of Crisis. \textit{Entropy}, 26(12), 1079.

\bibitem[Politis \& Romano(1994)]{politis1994} Politis, D. N., \& Romano, J. P. (1994). The Stationary Bootstrap. \textit{Journal of the American Statistical Association}, 89(428), 1303--1313.

\bibitem[Tauchen \& Pitts(1983)]{tauchen1983} Tauchen, G. E., \& Pitts, M. (1983). The Price Variability--Volume Relationship on Speculative Markets. \textit{Econometrica}, 51(2), 485--505.

\bibitem[Vasicek(1976)]{vasicek1976} Vasicek, O. (1976). A Test for Normality Based on Sample Entropy. \textit{Journal of the Royal Statistical Society B}, 38(1), 54--59.

\bibitem[Zhou, Cai \& Tong(2013)]{zhou2013} Zhou, R., Cai, R., \& Tong, G. (2013). Applications of Entropy in Finance: A Review. \textit{Entropy}, 15(11), 4909--4931.

\end{thebibliography}

\end{document}
